\section{Related Work}

\subsection{LLM-Based Test and Scenario Generation}
Recent advancements in large language models (LLMs) have driven significant research into automated software testing and scenario generation. Various models, including GPT-3, GPT-4, and CodeBERT, have been evaluated for their ability to synthesize test cases and behavior-driven development (BDD) scenarios from natural language artifacts \cite{wang2024llm, yuan2023chattester, siddiq2023exploring}. For instance, Wang et al. demonstrated the efficacy of ChatGPT in generating unit tests, albeit noting challenges in producing compilable code without human intervention \cite{wang2024llm}. Similarly, approaches leveraging CodeBERT have shown promise in mapping natural language requirements to code snippets, though they often struggle with the rigid structural constraints of specific domain-specific languages like Gherkin \cite{bareiss2022code}.

Evaluations in this domain have primarily relied on metrics such as BLEU scores, exact string matching, pass rates, and extensive human evaluation. While these metrics provide baseline indicators of syntactical similarity or functional correctness, they frequently fail to capture the semantic nuances of BDD scenarios. A shared limitation across these studies is their heavy reliance on qualitative human evaluation due to the lack of robust, automated quantitative metrics tailored for scenario generation. In a closely related study, Rathnayake et al. evaluated GPT-4 for zero-shot Gherkin generation, finding that it scored an average of 4.63 out of 5.0 based on assessments by human experts \cite{rathnayake2026llm}. While these results highlight the impressive capabilities of advanced LLMs, the dependence on manual scoring limits the scalability and objective reproducibility of such evaluations.

\subsection{BDD and Gherkin in Software Testing}
Behavior-Driven Development (BDD) has become a cornerstone of agile software development by fostering collaboration between technical and non-technical stakeholders through a shared ubiquitous language, most commonly Gherkin \cite{smart2014bdd}. Studies have shown that integrating BDD reduces miscommunication and aligns software behavior more closely with business requirements \cite{solis2011study}. However, the manual authoring of Gherkin scenarios is a widely recognized challenge. Research indicates that crafting and maintaining these scenarios can consume between 30\% and 50\% of the overall testing effort, posing a significant bottleneck in fast-paced agile environments \cite{irshad2021systematic}.

To address this, several automated approaches have been proposed. Early attempts focused heavily on rule-based natural language processing and template-driven generation, where user stories were parsed into predefined Gherkin structures using part-of-speech tagging and syntactic trees \cite{wang2021automated, alfadel2021empirical}. While effective for highly constrained inputs, these traditional automation methods proved brittle and struggled to generalize across diverse, unstructured user stories commonly found in industrial settings. Most importantly, none of these prior approaches---whether rule-based or early AI-driven---have validated their generation quality using robust, automated quantitative metrics on a large scale, leaving a significant gap in the objective assessment of automated BDD authoring tools.

\subsection*{}
Unlike prior work, which relies on human scoring or BLEU-based metrics, this paper is the first to evaluate GPT-4o zero-shot Gherkin generation using cosine semantic similarity with a fixed threshold ($\geq 0.85$) and executable syntax rate, validated with Wilcoxon and Binomial significance tests on an industrial dataset.
